La siniestralidad vial continúa siendo un problema de primer orden en las ciudades colombianas. Según el informe mundial de la Organización Mundial de la Salud, cada año mueren más de un millón de personas en accidentes de tránsito y millones más quedan con secuelas permanentes, con una carga desproporcionada en países de ingresos medios \cite{oms2024}. En Medellín, los registros de accidentalidad publicados en la plataforma de datos abiertos Mede ofrecen detalle sobre ubicación, hora, tipo de evento y gravedad de las víctimas; sin embargo, esa información suele consultarse de forma fragmentada y sin herramientas que integren mapa, indicadores y proyecciones en un mismo entorno \cite{mina2024,parraga2024}.\\
Este trabajo de grado, titulado \textit{Sistema de mitigación de accidentes de tránsito para ciudad con grandes urbes: caso Medellín}, presenta el diseño e implementación de \textbf{ViaData — Medellín}, plataforma web desarrollada para explorar esos datos históricos (periodo aproximado 2014--2021) y apoyar el análisis de la accidentalidad en el territorio urbano. La solución integra un backend en Django con PostgreSQL y extensión PostGIS, y un frontend en React que expone, de manera pública, un tablero de indicadores y un mapa analítico con capas de puntos, densidad y coroplética territorial; para perfiles analista y administrador, incorpora además un módulo de predicciones, reportes imprimibles y un asistente conversacional que consulta las mismas rutas de la API del sistema, sin generar cifras por cuenta propia.\\
El procesamiento de datos sigue un flujo ETL reproducible a partir del conjunto Mede, con reglas de depuración centralizadas y carga a base relacional georreferenciada. Sobre esa base, el módulo predictivo compara modelos estadísticos transparentes ---regresión lineal, estacionalidad, Poisson, media móvil, bandas $\mu \pm 3\sigma$, ARIMA y SARIMA--- mediante validación \textit{hold-out} sobre los últimos tres meses del rango analizado. Para la proyección mensual de incidentes a nivel ciudad, el modelo adoptado fue SARIMA(2,1,3)(1,1,1,12), con un error porcentual absoluto medio en prueba del 12,6\,\%. El sistema complementa esa proyección con un índice de prioridad territorial, estimación de carga esperada por comuna, evolución de la proporción de víctimas fatales y patrones de concentración por día de la semana y franja horaria.\\
Las proyecciones se plantean como escenarios orientativos bajo estabilidad del patrón histórico filtrado: no pretenden establecer causalidad ni estimar probabilidad individual de sufrir un siniestro. En conjunto, ViaData — Medellín materializa un instrumento geoespacial que conecta datos abiertos, visualización interactiva y analítica predictiva documentada, con miras a fortalecer la lectura territorial del riesgo vial en contextos metropolitanos \cite{monar2025}.

\palabrasclave{ViaData, accidentalidad vial, datos abiertos Mede, sistemas de información geográfica, series temporales, seguridad vial urbana.}
